\section{The tabs}

At desktop width every tab's panel is visible at once, stacked on one
scrolling page; the tab row itself only appears at phone width, where it
switches which single panel is shown. What follows is every panel, in the
order they appear on the page.

\subsection{Organization key tree}

The org-wide split tree, drawn as the dotted-label hierarchy it is:
\code{M}, \code{M.S}, \code{M.S.1}, and so on. You only ever see your own
slice --- your lineage, your descendants, siblings, and established-bridge
peers. A leaf tagged \emph{required} contributed a share the last time you
unlocked a file; \emph{satisfied} means that share was actually presented.

\begin{itemize}
  \item \textbf{Revoke key} bans a leaf's hardware key outright; it drops
  that leaf's bindings and pairings but does not evict it or refresh
  survivor shares.
  \item \textbf{Reissue\ldots} replaces a leaf's hardware key onto an
  already-provisioned replacement token.
  \item \textbf{Bridges} whitelists and establishes cross-branch links, the
  same two-step allow-then-add process the CLI uses.
  \item \textbf{Tree restructure} walks through proposing a new public
  generation and collecting the parent's countersignature when one is
  required.
\end{itemize}

\Screenshot{docs/latex/images/lab/organization.png}{The organization key
tree, showing the dotted-label hierarchy, per-leaf slot status, and the
Revoke/Reissue controls.}

\subsection{Mock USB devices}

Each card is one device container: a signed device descriptor plus one
sealed token per slot. \textbf{Insert} / \textbf{Eject} plug a drive in or
pull it out. \textbf{Move} asks for a destination drive (both drives must
be inserted) and relocates a slot's keypair onto it --- the point where two
people's tokens land in one physical container and start counting as a
single device toward any multi-device quorum. \textbf{Copy} instead leaves
the original slot active and seals a second copy elsewhere, using the same
passphrase that already unlocks it --- a seeded slot's own published default
passphrase (Section~\ref{sec:lab-default-credentials}).

\Screenshot{docs/latex/images/lab/usb.png}{Mock USB devices: connected
state, the slots each drive carries, and the Move / Copy controls.}

\subsection{File Explorer}

A Windows-Explorer-style view onto the active user's files, grouped into
folders. Double-clicking a file (or selecting it and pressing
\textbf{Open}) unlocks and views it --- a quorum file prompts for the
shares it needs, a password file for its password. \textbf{Send\ldots}
seals a copy to another label's key and drops it in their inbox;
\textbf{Sign} is only offered on public or received plaintext;
\textbf{Properties} shows the file's protection, size, and timestamps.
Right-clicking a row offers the same actions. A file with an expiry stamp
shows a status pill that flips to \emph{expired} on its own while the tab
is open.

\Screenshot{docs/latex/images/lab/files.png}{The File Explorer, showing the
public folder's contents and their status pills.}

\subsection{Inbox and sent}

Letters are sealed to a recipient's key and routed by the mailbox relay
without it ever reading them --- sender and contents stay hidden until
\textbf{Receive} unseals a letter with your own slot. \textbf{Reject}
answers with a signed refusal instead. The Sent view tracks whether a
delivery is still awaiting the recipient's signed acknowledgement, and
\textbf{Check relay for acknowledgements} pulls any that are waiting.

\Screenshot{docs/latex/images/lab/mailbox.png}{Inbox and sent, in the empty
seeded state.}

\subsection{Security \& devices}

Everything that is not the hardware-key quorum tree or a file: locking a
note with a password (and optionally a PIN), exporting a file or credential
as a portable sealed bundle, minting a time-limited share link,
provisioning a fresh slot on an inserted drive and registering it as a new
leaf, signing or verifying with a private sign bridge you belong to, and,
further down, read-only device logs and relay counters for the in-process
mailbox this lab runs against.

\Screenshot{docs/latex/images/lab/security.png}{Security \& devices: the
password-lock form, exported bundles, share links, and key/leaf
provisioning.}

\subsection{Activity and access trace}

Every action taken leaves a trace here: which checks ran, in order, and
whether each one passed. The \textbf{Equivalent operation} line under the
latest trace is the literal command line that action ran --- the same one
that could be typed in the Terminal tab. The full log below keeps every
action for the session, not just the last one.

\paragraph{Tracked files.} The Activity tab is also where tracked files
live. The \textbf{Tracked files} panel acts as whoever is active, with their
own store and their own drive's slot; each button runs one real
\code{keyquorum file} command:

\begin{itemize}
  \item \textbf{Track and sign} a new file, then \textbf{Check in} edits ---
  signed with your slot, or unsigned to see a revision stay
  \emph{pending}.
  \item Per revision: \textbf{Sign revision} (its author),
  \textbf{Approve as parent} (the author's direct parent countersigning),
  \textbf{Finalize} (the scope owner or an ancestor, once it is trusted),
  \textbf{View} an older revision, and \textbf{Diff} it against its parent.
  Each revision shows its user label first, when it has one, then its
  generated label and short id.
  \item \textbf{Rename file}: change the name shown; the file id and every
  revision id stay the same.
  \item \textbf{Share file} with another person. The line above the button
  says which revision would go --- the head, or the last trusted revision
  when the head is pending. The recipient \textbf{Accept}s or
  \textbf{Refuse}s the letter (into their own copy), and the sender
  \textbf{Record}s \textbf{the answer}.
  \item \textbf{Ask for a file, or for a change}: pick a person, choose
  \emph{a file} (send me your copy) or \emph{a change to this file}, add a
  message if you like, and press \textbf{Send request} (\code{keyquorum file
  request}, with \code{--change} for a change). Under \emph{Requests for a
  file or a change} the holder sees the message and presses \textbf{Accept
  request} or \textbf{Decline request} (\code{open-request}, then
  \code{answer-request}, which also records the request in their own copy
  of the file). The asker presses \textbf{Record the request answer}
  (\code{open-answer}). A request only asks: nothing is delivered or
  changed. To serve an accepted file request the holder uses \textbf{Share
  file}; a change arrives as an ordinary revision. The log shows the request
  id, kind and decision, never the message.
  \item When two copies diverge: \textbf{Import copy}, \textbf{Review the
  fork} (each side's changed lines and the named reviewer) and
  \textbf{Merge heads}; a clean merge is a new revision that still has to be
  signed. \textbf{Resolve} runs \code{keyquorum file resolve} as the active
  person with their own slot: keep the left or right side, use text you
  type, or reject a proposed merge. Only the assigned reviewer (or an
  ancestor of theirs) can settle a conflict, so the same button is refused
  for anyone else; the lab adds no check of its own.
  \item \textbf{Verify history}, \textbf{Export history snapshot} and
  \textbf{Check} a snapshot against the file.
  \item \textbf{Schedule expiry} or \textbf{Destroy content now}: the file
  becomes a tombstone that still verifies.
  \item \textbf{Link gate}: record unlocks of a quorum file, or of one of
  your password files, in this file's history; \textbf{Unlink} stops it.
\end{itemize}

Whatever a button or a Terminal command writes to a tracked file appears in
the log at once, right below the action that caused it. History entries can
be expanded (revision, labels, parents, trust, history root) and filtered
by \textbf{Tracked File}, \textbf{Revision}, \textbf{Security},
\textbf{Sharing} and \textbf{Conflict}; a revision graph per file sits
under the log. Three tracked files are seeded: \file{budget.txt} (an
unsigned edit, so sharing falls back), \file{forecast.txt} (two edits that
merged on their own) and \file{memo.txt} (two edits to one line, assigned
to a reviewer).

\Screenshot{docs/latex/images/lab/activity.png}{The Activity tab's tracked
files panel: revisions with their trust, what sharing would send, and the
controls for the selected file.}

\subsection{Terminal}

A real shell on the same lab machine. Every button on every other tab runs
a command here under the hood; typing it directly works identically,
including commands no button exposes yet.
Use full paths in \code{keyquorum} arguments: the lab's shell does not
expand \code{~} inside them, so \flag{--db ~/keyquorum.sqlite} would open a
new, empty store rather than your own.

\begin{lstlisting}
keyquorum-device list /media/alice-usb
bridge list 1
usb insert david
su david
\end{lstlisting}

\code{usb insert} / \code{usb eject} and \code{su} are lab-only shell verbs
for drive and identity control; every other line is an ordinary
\code{keyquorum} / \code{keyquorum-device} command. Any passphrase,
password, or PIN prompt a typed command raises is answered with its
published default value (Section~\ref{sec:lab-default-credentials}), the
same way a GUI action's prompt is when you have not staged your own
answer first.
